\section{Discussion}
\label{sec:discussion}

\paragraph{Limitations.}
All closed-loop evaluations use the LIBERO simulation benchmark; real-robot validation is absent. The six models differ in backbone capacity, vision encoder, and fusion strategy beyond the connector, so the connector effect cannot be causally isolated. $\pi_0$ is evaluated through offline metrics only, without closed-loop verification. The observed correlation between CV and pruning tolerance spans only three closed-loop models ($\pi_0.5$, X-VLA, OpenVLA), too few to establish a general relationship. SmolVLA baselines (59.2\% official, 44.8\% LoRA reproduction) are substantially weaker than those of the other models; whether VTR recovery generalizes to stronger baselines is unknown.

\paragraph{Non-monotonicity as structural property.}
X-VLA on Goal exhibits non-monotonic degradation: $k{=}0.95$ drops to 53.8\% (L2) and 54.0\% (random), while $k{=}0.9$ recovers to 70.2\% under L2 selection (Table~\ref{tab:main_results}). The 0.2\,pp method gap at $k{=}0.95$ confirms that the anomaly is independent of the token scoring criterion. This pattern is consistent with position sensitivity in the 12-layer BART encoder~\cite{lewis2020bart}, where specific token positions carry disproportionate structural importance regardless of content.

\paragraph{$\pi_0.5$ random subsampling anomaly.}
Random selection on $\pi_0.5$ Object collapses from 69.4\% at $k{=}0.95$ to 12.0\% at $k{=}0.9$, then recovers to 77.2\% at $k{=}0.7$ (Table~\ref{tab:main_results}). The V-shaped recovery at more aggressive pruning remains unexplained.

\paragraph{Small-$N$ interactions.}
X-VLA operates with only 50 vision tokens. At $k{=}0.95$, removing 3 tokens produces $-$44.2\,pp on Goal and $-$47.0\,pp on Long. Absolute token count interacts with CV: identical retention ratios impose qualitatively different damage when the token budget is small.

\paragraph{Design implications.}
Practitioners should compute importance CV before deploying token pruning, as low-CV and entangling-connector models require distinct compression strategies. Pixel-shuffle connectors need learned re-encoding rather than hard deletion. Task complexity must factor into pruning budget decisions: retention ratios safe for simple manipulation may fail on compositional tasks.